\documentclass[10pt,a4paper]{amsart}
\usepackage[margin=19mm]{geometry}
\usepackage{amsmath,amssymb,mathtools}
\usepackage[scr=rsfs,cal=euler]{mathalfa}
\usepackage{xfrac}
\usepackage[unicode,hidelinks,bookmarks=false]{hyperref}
\newcommand{\ie}{i.e.\ }
\newcommand{\cf}{cf.\ }
\newcommand{\resp}{respectively\ }
\newcommand{\Subsection}{\subsection}
\providecommand{\nobreakdash}{\nobreak\hbox{-}}
\setlength{\emergencystretch}{2em}
\multlinegap0pt
\usepackage{bm}
\usepackage{bbm}
\usepackage{dsfont}
\usepackage{xspace}
\usepackage{cancel}
\usepackage{mathtools}
\usepackage{amsthm}
\makeatletter
\@ifundefined{Lemma}{\newtheorem{Lemma}{Lemma}}{}
\@ifundefined{Proposition}{\newtheorem{Proposition}{Proposition}}{}
\makeatother
\newcommand{\Dv}{{\rm div}}
\newcommand{\R}{{\mathbb R}}
\newcommand{\id}{\mathrm{Id}}
\newcommand{\trace}{\mathrm{tr}}
\newcommand{\tr}{{\rm tr}}
\title[Universality in the Low Mach number limit via a convex integ]{Universality in the Low Mach number limit via a convex integration framework}
\author{Robin Ming Chen, Alexis Vasseur, Dehua Wang, Cheng Yu}
\date{Source: arXiv:2601.19744v1 (CC BY 4.0)}
\begin{document}
\maketitle

\noindent\emph{Source and licence.} Universality in the Low Mach number limit via a convex integration framework. Version arXiv:2601.19744v1 (CC BY 4.0), distributed under the Creative Commons Attribution 4.0 International licence, as recorded in the arXiv licence metadata for this exact version and shown on the abstract page https://arxiv.org/abs/2601.19744v1. Full rights notice: licenses/arxiv-2601.19744-CC-BY-4.0.txt.\par\smallskip

\noindent\textit{Editorial scope.} Counted unit: the Proposition whose source label is \texttt{proposition points} in the article (source0.tex lines 694--700, source label \texttt{proposition points}), delivered together with the article's own complete original proof of it (source0.tex lines 701--757, one \texttt{proof} environment of 57 lines). No mathematics is written, repaired or re-derived: statement and proof are the article's own text line for line. Required context retained uncounted: ci system (lines 434--441); restriction = (lines 443--446); bounds rho (lines 464--467); additional control (lines 471--474); ci constraint 2 (lines 480--483); defn X\_0 (lines 485--487); L1-coercivity (lines 502--509). Every definition, display and label of those lines is retained unchanged. Omitted from this package and not redistributed: 1--433, 442--442, 447--463, 468--470, 475--479, 484--484, 488--501, 510--693, 758--977. Nothing inside a retained excerpt is omitted or altered: every retained body is delivered exactly as the article's lines are written, so no omission receipt of any kind accompanies this package. Citations: the retained excerpts carry no citation command at all, so no bibliography entry of the article is transcribed and no reference is redistributed. Reference adaptation: none was needed and none was made; every reference inside the delivered unit resolves inside it, so no unresolved reference or citation is produced. Printed slips kept literal: no printed word, symbol, hypothesis or number of the counted statement or of its proof was altered. Layout and class: the article's own class is \texttt{amsart} with packages including graphicx, fullpage, bbm, amssymb, amsmath, amsthm, tikz, bm, epstopdf, enumitem, esint, verbatim, mathabx, hyperref; the wrapper is the lane's pinned \texttt{amsart} editorial wrapper at 10pt on A4, which restates the theorem-like environments the retained text needs and adds the article's own macro definitions that the retained text needs (\texttt{\textbackslash Dv}, \texttt{\textbackslash R}, \texttt{\textbackslash id}, \texttt{\textbackslash tr}, \texttt{\textbackslash trace}), with extra wrapper packages (bm, bbm, dsfont, xspace, cancel, mathtools). The delivered numbering is the unit's own. Assets: no retained span contains a figure environment, a graphics-inclusion command, a PDF-page inclusion, a table or a TikZ picture; no image file is copied into this package, the package declares no asset dependency and no image file is redistributed with it. Licence: version arXiv:2601.19744v1 of this article is distributed under the Creative Commons Attribution 4.0 International licence, as recorded in the arXiv licence metadata for this exact version and shown on the abstract page https://arxiv.org/abs/2601.19744v1. A full rights notice is delivered at licenses/arxiv-2601.19744-CC-BY-4.0.txt. Characters: every retained excerpt is pure ASCII, so the delivered engine, XeLaTeX with its default Latin Modern text font, reports no missing character.
\begin{equation}
\label{ci system}
 \left\{
\begin{split}
&\Dv \widetilde{V}=0,\\
&\partial_t\widetilde{V}+\Dv \widetilde{U}=0,
\end{split}\right.
\end{equation} 

\begin{equation}
\label{restriction =}
\frac{(V_0 + \widetilde{V}) \otimes (V_0 + \widetilde{V})}{\rho_0} - (U_0 + \widetilde{U}) = \frac{|V_0|^2}{n\rho_0} \id_n + R_0 \quad \text{a.e. in } P,
\end{equation}

\begin{equation}
 \label{bounds rho}
 0<\frac{1}{\Lambda^2}\leq \rho_0\leq \Lambda^2.
 \end{equation}

\begin{equation}
\label{additional control}
\left|\int_0^T\int_{\mathbb{T}^3}\frac{V_0\cdot\widetilde{V}}{\rho_0}\,dx\,dt\right|\leq \frac{1}{8}\int_0^T\int_{\mathbb{T}^3} \frac{|\widetilde{V}|^2}{\rho_0}\,dx\,dt.
\end{equation}

\begin{equation}
\label{ci constraint 2}
\frac{(V_0 + \widetilde{V}) \otimes (V_0 + \widetilde{V})} { \rho} - (U_0 + \widetilde U) < {|V_0|^2 \over n\rho_0} \id_n + R_0,
\end{equation}

\begin{equation}\label{defn X_0}
X_0 := \left\{ (\widetilde V, \widetilde U) \in C^\infty_c(P; \R^n \times \mathbb{S}^{n\times n}_0): \ (\widetilde V, \widetilde U) \text{ solves } \eqref{ci system},   \eqref{additional control} \text{ and } \eqref{ci constraint 2} \right\},
\end{equation}

\begin{Lemma}[$L^1$-coercivity]
\label{L1-coercivity}
For any $(\widetilde{V}, \widetilde{U}) \in X_0$, where $X_0$ is defined in \eqref{defn X_0}, there exists a sequence $\{(\widetilde{V}_i, \widetilde{U}_i)\} \subset X_0$ that converges weakly-$*$ to $(\widetilde{V}, \widetilde{U})$, such that
$$
\|\widetilde{V}_i - \widetilde{V}\|_{L^1(P)} \geq \frac{c_0}{\Lambda} \int_{P} \trace \, M \, dx \, dt,
$$
for some positive geometric constant $c_0$, where $\Lambda$ is given in \eqref{bounds rho}.
\end{Lemma}

\begin{Proposition}
\label{proposition points}
Let $(\widehat{V},\widehat{U})\in X$ be a point of continuity of the identity map $\id$ from $(X,d)$ to $L^2(\R^n\times\R)$. Then, $(\widehat{V},\widehat{U})$ satisfies \eqref{restriction =}. In addition, there exists a positive number $C_0>0$, such that 
 $$
 \int_0^T\int_{\mathbb{T}^n}\frac{|\widehat{V}|^2}{\rho_0}\,dx\,dt\leq C_0\int_0^T\int_{\mathbb{T}^n} \trace R_0\,dx\,dt.
 $$
\end{Proposition}

\begin{proof}
By definition, there exists a sequence $\{(V_j, U_j)\} \subset X_0$ that converges weak-* to $(\widehat{V}, \widehat{U})$. In particular, we have
\[
V_j \to \widehat{V} \quad \text{strongly in } L^2(P),
\]
and hence also strongly in $L^1_{\text{loc}}(P)$.

From Lemma \ref{L1-coercivity}, for each $(V_j, U_j)$ there exists a subsequence $\{(V_{j,i}, U_{j,i})\}_{i\in \mathbb{N}}$ converging weak-* to $(V_j, U_j)$. Note that $(V_j, U_j)$ satisfies \eqref{ci system}, and we have the estimate
\[
\|V_{j,i} - V_j\|_{L^1(P)} \geq \frac{c_1}{8\Lambda} \int_P \trace M_j \, dx\,dt,
\]
where $M_j$ is given by
\[
M_j := \frac{|V_0|^2}{n\rho_0} \id_n + R_0 - \frac{(V_0 + V_j) \otimes (V_0 + V_j)}{\rho_0} + U_0 + U_j.
\]

By a diagonal argument, we may extract a subsequence $\{(V_{j,i(j)}, U_{j,i(j)})\}$ that converges weakly-* to $(\widehat{V}, \widehat{U})$ and satisfies
\[
\lim_{j \to \infty} \|V_{j,i(j)} - \widehat{V}\|_{L^1(P)} \geq \frac{c_1}{8\Lambda} \int_P \trace M_j \, dx\,dt.
\]
  This yields to 
\begin{equation}
  \label{limit of Mj}
  \lim_{j}\int_{P}\trace M_j\,dx\,dt=0.
\end{equation}
  
We further introduce  
\begin{equation*}
  \widehat{M} := \frac{|V_0|^2}{n\rho_0} \id_n +R_0-\frac{(V_0+\widehat{V})\otimes (V+\widehat{V})}{\rho_0}+U_0+\widehat{U}.
  \end{equation*}
It follows that $M_j\to \widehat{M}$ weak-$*$. Note that $M_j>0$, which implies that $\widehat{M}\geq 0$ a.e.. From \eqref{limit of Mj} and noting that 
$$
\trace M_j\to \trace \widehat{M}\;\;\text{weak-}*,
$$
one obtains  
\begin{equation*}
  \lim_{j}\int_{P} \trace M_j\,dx\,dt = \int_P\tr \widehat{M}\,dx\,dt=0.
\end{equation*}
   Thus, $\tr \widehat{M}=0$ and so $\widehat{M}=0$, which yields that $(\widehat{V},\widehat{U})$ satisfies  \eqref{restriction =}.
    
Meanwhile, note from the definition of $\widehat{M}$ that 
  \begin{equation*}
  \begin{split}
  \tr \widehat{M}&=\frac{|V_0|^2}{\rho_0} + \tr R_0-\frac{|V_0+\widehat{V}|^2}{\rho_0}
  = \tr R_0-\frac{|\widehat{V}|^2}{\rho_0}-\frac{2V_0\cdot \widehat{V}}{\rho_0}.
  \end{split}
  \end{equation*}  
  This yields that
  \begin{equation*}
  \begin{split}
  \int_P\tr R_0\,dx&=\int_P\frac{|\widehat{V}|^2}{\rho_0}\,dx+2\int_P\frac{V_0\cdot\widehat{V}}{\rho_0}\,dx
  \\&\geq \int_P\frac{|\widehat{V}|^2}{\rho_0}\,dx-\frac{1}{4}\int_P\frac{|\widehat{V}|^2}{\rho_0}\,dx
  \\&=\frac{3}{4}\int_P\frac{|\widehat{V}|^2}{\rho_0}\,dx,
  \end{split}
  \end{equation*}
completing the proof.   
\end{proof}

\end{document}
